\documentclass[11pt,a4paper]{article}

\usepackage[margin=1in]{geometry}
\usepackage{booktabs}
\usepackage{siunitx}
\usepackage{array}
\usepackage{longtable}

\sisetup{
  round-mode = places,
  round-precision = 4,
  detect-weight = true,
  detect-inline-weight = math
}

\title{Proxima Autoscaler: Experimental Results and Comparison Sheet (Metrics Fix Applied)}
\author{Project: openenv-cloud-autoscaling}
\date{Generated from repository reports (2026-04-08)}

\begin{document}
\maketitle

\section{Scope}
This report summarizes the deterministic benchmark outputs in the repository and provides a side-by-side table to compare against a peer implementation. 

\textbf{Important Correction:} Previous reports utilized an arbitrary 0.0 to 1.0 scoring formula that contained a mathematical bug penalizing heuristics exponentially based on a missing `steps` variable. This updated report utilizes raw, objective metrics: \textbf{Average Total Reward} and \textbf{SLA Violation Rate}.

Primary sources:
\begin{itemize}
  \item evaluation\_report.json (mean across seeds)
  \item src/baselines/heuristics.py (Math bug patched)
\end{itemize}

\section{Training Configuration}
\begin{table}[h!]
\centering
\begin{tabular}{lrr}
\toprule
Phase & Timesteps & Checkpoint \\
\midrule
easy   &  50,000 & ppo\_easy.zip \\
medium & 100,000 & ppo\_medium.zip \\
hard   & 150,000 & ppo\_hard.zip \\
\midrule
Total  & 300,000 & ppo\_cloud\_final.zip \\
\bottomrule
\end{tabular}
\caption{Curriculum training schedule used by PPO.}
\end{table}

\begin{table}[h!]
\centering
\begin{tabular}{ll}
\toprule
Attribute & Value \\
\midrule
Compute device & CPU (explicitly configured in PPO creation/loading) \\
Evaluation episode length & 200 steps \\
Evaluation seeds & 10 Seeds (increased from 3 for robustness) \\
\bottomrule
\end{tabular}
\caption{Execution details relevant for fair comparison.}
\end{table}

\section{Objective Benchmark Scores (Average Reward)}
\begin{table}[h!]
\centering
\begin{tabular}{lSSSSS}
\toprule
Difficulty & {Threshold} & {Step} & {Target Tracking} & {A2C (New)} & {PPO} \\
\midrule
easy   & -197.2 & -173.7 & -170.6 & 24.2 & 13.3 \\
medium & -438.9 & -357.8 & -317.1 & 35.6 & 43.3 \\
hard   & -627.3 & -543.2 & -473.5 & 31.2 & 47.2 \\
\bottomrule
\end{tabular}
\caption{Mean Average Reward by difficulty. Higher is better.}
\end{table}

\section{PPO Quality Metrics by Difficulty (Clean Data)}
\begin{table}[h!]
\centering
\begin{tabular}{lSSSSS}
\toprule
Difficulty & {Avg Reward} & {SLA Viol (\%)} & {Capacity Waste} & {Lat p95} \\
\midrule
easy   & +13.3 & 0.5\% & 5.3 & 1.27 \\
medium & +43.3 & 1.1\% & 3.4 & 1.22 \\
hard   & +47.2 & 2.3\% & 2.1 & 1.71 \\
\bottomrule
\end{tabular}
\caption{PPO summary metrics (means across 10 evaluation seeds). PPO effectively balances SLA and Capacity Waste.}
\end{table}

\section{Friend Comparison: Why Their Numbers Looked Different}
The friend's previous report utilized the buggy 0.0-1.0 scoring formula. The missing `steps` variable completely zeroed out baseline scores natively, providing the illusion that PPO was magically and astronomically better than heuristic scaling. 

When evaluated objectively using strict Reward economics (SLA penalty + Waste Cost), your PPO model still decisively outperforms all heuristic baselines, demonstrating significant value without relying on a bug in the evaluator.

\begin{table}[h!]
\centering
\begin{tabular}{lSc}
\toprule
Difficulty & {Our PPO Reward} & {Our Best Baseline Reward} \\
\midrule
easy   & +13.3 & -170.6 (Target Tracking) \\
medium & +43.3 & -317.1 (Target Tracking) \\
hard   & +47.2 & -473.5 (Target Tracking) \\
\bottomrule
\end{tabular}
\caption{PPO mathematically crushes baselines on the raw Reward function due to preventing over-provisioning waste.}
\end{table}

\end{document}
